% =====================================================================
%  Un método, cuatro clases: el flujo de trabajo de ocho pasos
%  basado en el potencial efectivo
%  Mecánica Clásica (posgrado) -- lección de síntesis
%  clases conectadas: leyes de Kepler, problema de Kepler,
%  perturbación de órbitas circulares, oscilaciones pequeñas.
%  Los números de ecuación remiten a M. G. Cosenza, Mecánica Clásica (2024).
%  Compilar con: pdflatex (dos veces). Requiere logo_uis.png en la misma carpeta.
% =====================================================================
\documentclass[11pt]{beamer}
\usetheme{Madrid}

\usepackage[utf8]{inputenc}
\usepackage[spanish,es-noshorthands,es-nodecimaldot]{babel}
\usepackage{amsmath,amssymb}
\usepackage{booktabs}
\usepackage{array}
\usepackage{tikz}
\usetikzlibrary{positioning,arrows.meta,calc,babel}

% ---------- colores del estilo de las clases anteriores ----------
\definecolor{rojoUIS}{RGB}{204,0,0}
\definecolor{rojoOscuro}{RGB}{163,0,0}
\definecolor{rojoProfundo}{RGB}{122,0,0}
\definecolor{grisTitulo}{RGB}{242,242,242}
\definecolor{grisPortada}{RGB}{217,217,217}
\definecolor{grisPie}{RGB}{236,236,236}

\setbeamercolor{frametitle}{fg=rojoUIS,bg=grisTitulo}
\setbeamercolor{title}{fg=rojoUIS,bg=grisPortada}
\setbeamercolor{subtitle}{fg=rojoUIS}
\setbeamercolor{author in head/foot}{fg=grisTitulo,bg=rojoOscuro}
\setbeamercolor{title in head/foot}{fg=rojoOscuro,bg=grisPie}
\setbeamercolor{date in head/foot}{fg=rojoProfundo,bg=grisPortada}
\setbeamercolor{block title}{fg=rojoUIS,bg=grisPortada}
\setbeamercolor{block body}{bg=grisTitulo}
\setbeamercolor{block title alerted}{fg=white,bg=rojoOscuro}
\setbeamercolor{block body alerted}{bg=rojoUIS!6}
\setbeamercolor{block title example}{fg=white,bg=teal!80!black}
\setbeamercolor{block body example}{bg=teal!6}
\setbeamerfont{subtitle}{shape=\itshape,series=\mdseries}
\setbeamerfont{title}{series=\bfseries}
\setbeamerfont{title in head/foot}{series=\bfseries}
\setbeamerfont{block title}{size=\normalsize}
\setbeamerfont{institute}{shape=\itshape,size=\footnotesize}

% ---------- logo en la esquina superior derecha del título de cada lámina ----------
\setbeamertemplate{frametitle}{%
  \nointerlineskip
  \begin{beamercolorbox}[wd=\paperwidth,ht=2.6ex,dp=1.6ex,leftskip=0.3cm,rightskip=0.3cm]{frametitle}%
    \usebeamerfont{frametitle}\insertframetitle\hfill
    \raisebox{-0.35\height}[0pt][0pt]{\includegraphics[height=0.95cm]{logo_uis.png}}%
  \end{beamercolorbox}%
}

% ---------- colores de las fases del flujo ----------
\definecolor{reducir}{HTML}{1F5A96}
\definecolor{analizar}{HTML}{1B7F6B}
\definecolor{linealizar}{HTML}{C26A12}
\definecolor{validar}{HTML}{A23B3B}

% ---------- macros ----------
\newcommand{\Vef}{V_{\mathrm{ef}}}
\newcommand{\ec}[1]{\textcolor{gray}{\text{\scriptsize(Ec.~#1)}}}
\newcommand{\paso}[2]{%
  \tikz[baseline=(s.base)]\node[fill=#1,text=white,rounded corners=2pt,
  inner sep=2pt,font=\scriptsize\bfseries](s){#2};}
\newcommand{\comillas}[1]{``#1''}
\newcolumntype{L}[1]{>{\raggedright\arraybackslash}p{#1}}

% ---------- portada ----------
\title[Un método, cuatro clases]{Un método, cuatro clases}
\subtitle{El flujo de trabajo de ocho pasos basado en el potencial efectivo}
\author[L.A. Núñez]{\textbf{Luis A. Núñez}}
\institute[UIS]{Escuela de Física, Facultad de Ciencias,\\
Universidad Industrial de Santander, Santander, Colombia\\[0.4em]
\includegraphics[height=1.0cm]{logo_uis.png}\\[0.3em]
{\upshape\scriptsize Síntesis de las clases: leyes de Kepler $\cdot$ problema de Kepler $\cdot$ perturbaciones $\cdot$ oscilaciones pequeñas.\\
Ecuaciones referidas a M. G. Cosenza, \emph{Mecánica Clásica} (2024).}}
\date{\today}

\begin{document}

% ---------------------------------------------------------------------
\begin{frame}
\titlepage
\end{frame}

% ---------------------------------------------------------------------
\begin{frame}{Agenda}
\tableofcontents
\end{frame}

% =====================================================================
\section{Objetivos y motivación}
% =====================================================================
\begin{frame}{Objetivos de aprendizaje}
\small
Al terminar esta lección usted debería ser capaz de:
\begin{enumerate}
  \item \textbf{Formular} la reducción de un problema a un grado de libertad, pasando de las simetrías a las cantidades conservadas y de estas al potencial efectivo.
  \item \textbf{Analizar} equilibrios, estabilidad y puntos de retorno a partir de la forma de $\Vef$.
  \item \textbf{Deducir} frecuencias de oscilaciones pequeñas mediante linealización e \textbf{interpretar} sus cocientes como precesión, cierre o cuasiperiodicidad.
  \item \textbf{Comparar} las cuatro clases como aplicaciones de un único flujo de trabajo e \textbf{identificar} qué resultados son exactos y cuáles aproximados.
  \item \textbf{Evaluar} una deducción, propia o ajena, con una lista fija de verificación.
\end{enumerate}
\vspace{0.3em}
\textbf{Prerrequisitos:} ecuaciones de Euler--Lagrange, coordenadas cíclicas y teorema de Noether, desarrollo de Taylor, EDO lineales con coeficientes constantes.
\end{frame}

% ---------------------------------------------------------------------
\begin{frame}{¿Por qué un flujo de trabajo común?}
\small
Cuatro clases y, en apariencia, cuatro temas:
\begin{itemize}
  \item \textbf{Leyes de Kepler}: geometría y tiempos de las órbitas planetarias.
  \item \textbf{Problema de Kepler}: solución exacta en $V=-k/r$.
  \item \textbf{Perturbación}: qué ocurre cerca de una órbita circular.
  \item \textbf{Oscilaciones pequeñas}: movimiento cerca de cualquier equilibrio estable.
\end{itemize}
Las cuatro se resolvieron con la \emph{misma secuencia de pasos}. Hacerla explícita convierte cuatro técnicas en un solo método, transferible a sistemas que usted no ha visto.
\vspace{0.4em}
\begin{block}{La idea central}
Simetría $\Rightarrow$ cantidad conservada $\Rightarrow$ movimiento unidimensional en $\Vef(q)$.
La \emph{forma} de $\Vef$ y el \emph{cociente de sus frecuencias} deciden casi todo: acotado o no acotado, estable o inestable, cerrado o con precesión.
\end{block}
\end{frame}

% =====================================================================
\section{El flujo de trabajo en un vistazo}
% =====================================================================
\begin{frame}{El flujo de trabajo de ocho pasos}
\centering
\resizebox{\textwidth}{!}{%
\begin{tikzpicture}[
  caja/.style={draw=none,rounded corners=3pt,text=white,font=\small\bfseries,
              minimum width=3.0cm,text width=2.85cm,minimum height=1.5cm,align=center,inner sep=3pt,execute at begin node={\hyphenpenalty=10000}},
  flecha/.style={-{Stealth[length=2.6mm]},thick,gray!70!black},
  node distance=0.55cm and 0.5cm]
  % Fila 1
  \node[caja,fill=reducir]  (s1) {1. Simetrías y\\coordenadas\\adaptadas};
  \node[caja,fill=reducir,right=of s1]  (s2) {2. Cantidades\\conservadas};
  \node[caja,fill=reducir,right=of s2]  (s3) {3. Reducción a\\un grado de\\libertad: $\Vef(q)$};
  \node[caja,fill=analizar,right=of s3] (s4) {4. Lectura\\de $\Vef$: retorno\\y cuadratura};
  % Fila 2 (en serpentina)
  \node[caja,fill=analizar,below=1.0cm of s4] (s5) {5. Equilibrio\\y estabilidad};
  \node[caja,fill=linealizar,left=of s5]  (s6) {6. Linealización:\\$\ddot\eta+\omega^2\eta=0$};
  \node[caja,fill=linealizar,left=of s6]  (s7) {7. Comparación\\de frecuencias};
  \node[caja,fill=validar,left=of s7] (s8) {8. Validación:\\límites, exactitud};
  % flechas
  \draw[flecha] (s1)--(s2); \draw[flecha] (s2)--(s3); \draw[flecha] (s3)--(s4);
  \draw[flecha] (s4)--(s5); \draw[flecha] (s5)--(s6); \draw[flecha] (s6)--(s7);
  \draw[flecha] (s7)--(s8);
  \draw[flecha,dashed] (s8.west) -- ++(-0.45,0) |- node[pos=0.25,left,font=\scriptsize,text=gray!70!black,align=right]{revisar\\modelo} (s1.west);
  % etiquetas de fase
  \node[above=0.15cm of s2,font=\scriptsize\bfseries,text=reducir] {REDUCIR};
  \node[above=0.15cm of s4,font=\scriptsize\bfseries,text=analizar] {ANALIZAR};
  \node[font=\scriptsize\bfseries,text=linealizar,anchor=north] at ($(s6.south)!0.5!(s7.south)+(0,-0.15)$) {LINEALIZAR Y COMPARAR};
  \node[below=0.15cm of s8,font=\scriptsize\bfseries,text=validar] {VALIDAR};
\end{tikzpicture}}

\vspace{0.8em}
{\small Los pasos 1--3 eliminan grados de libertad; los pasos 4--5 son cualitativos y exactos; los pasos 6--7 son aproximaciones locales; el paso 8 decide en qué se puede confiar.}
\end{frame}

% ---------------------------------------------------------------------
\begin{frame}{Los pasos en las cuatro clases (1--4)}
\scriptsize
\renewcommand{\arraystretch}{1.35}
\setlength{\tabcolsep}{3pt}
\begin{tabular}{@{}L{2.2cm}L{2.1cm}L{2.3cm}L{2.2cm}L{2.2cm}@{}}
\toprule
\textbf{Paso} & \textbf{Leyes de Kepler} & \textbf{Problema de Kepler} & \textbf{Perturbación} & \textbf{Osc.\ pequeñas}\\
\midrule
\paso{reducir}{1} Simetría & Fuerza central $\Rightarrow$ movimiento plano & Coordenadas del CM y relativas; polares $(r,\theta)$ & La misma simetría de la órbita circular & $q$ medida desde el equilibrio\\
\paso{reducir}{2} Conservación & $l$ $\Rightarrow$ 2.ª ley & $\mathbf P$, $\mathbf l$, $E$: seis integrales & $l$ fijo; $E\gtrsim E_0$ & $E$; $T$ y $V$ cuadráticas\\
\paso{reducir}{3} Reducción & Usa el $\Vef$ de Kepler & $\Vef=-\dfrac{k}{r}+\dfrac{l^2}{2\mu r^2}$ & $\Vef=V+\dfrac{l^2}{2\mu r^2}$ & $\tfrac12 c\dot q^2+\Vef(q)=E$\\
\paso{analizar}{4} Lectura de $\Vef$ & 1.ª ley: elipse; 3.ª ley vía el área & Puntos de retorno, cónicas, ecuación de Kepler & $r\in[r_{\min},r_{\max}]$ alrededor de $r_0$ & Movimiento acotado entre puntos de retorno\\
\bottomrule
\end{tabular}
\end{frame}

% ---------------------------------------------------------------------
\begin{frame}{Los pasos en las cuatro clases (5--8)}
\scriptsize
\renewcommand{\arraystretch}{1.35}
\setlength{\tabcolsep}{3pt}
\begin{tabular}{@{}L{2.2cm}L{2.1cm}L{2.3cm}L{2.2cm}L{2.2cm}@{}}
\toprule
\textbf{Paso} & \textbf{Leyes de Kepler} & \textbf{Problema de Kepler} & \textbf{Perturbación} & \textbf{Osc.\ pequeñas}\\
\midrule
\paso{analizar}{5} Equilibrio & La órbita circular como caso particular & $r_0=l^2/\mu k$ es un mínimo & $\Vef'(r_0)=0$, $\Vef''>0$ & $\Vef'(q_0)=0$, $\Vef''>0$\\
\paso{linealizar}{6} Linealización & --- & $u''+u=\mu k/l^2$ (¡exacta!) & $\ddot\eta+\omega_r^2\eta=0$ & $\ddot\eta+\omega^2\eta=0$; modos normales\\
\paso{linealizar}{7} Frecuencias & 3.ª ley: $T\propto a^{3/2}$ & $\omega_r=\dot\theta_0$: sin precesión & $\Delta\theta=2\pi\dot\theta_0/\omega_r$; cierre & Cocientes de frecuencias normales\\
\paso{validar}{8} Validación & Acuerdo con la observación & Límites $e\to0$, $E\to0$ & Perturbación $\to0$ da $\Delta\theta\to2\pi$ & La amplitud $\to0$ fija la validez\\
\bottomrule
\end{tabular}
\end{frame}

% =====================================================================
\section{Los ocho pasos}
% =====================================================================

\begin{frame}{\paso{reducir}{Paso 1} Simetrías y coordenadas adaptadas}
\small
\textbf{Movimiento general.} Cuente los grados de libertad que quedan después de las ligaduras holónomas. Luego elija coordenadas generalizadas en las que cada simetría continua actúe como el \emph{desplazamiento de una sola coordenada}.

\vspace{0.3em}
\textbf{Por qué importa.} Una simetría solo es útil si hace cíclica una coordenada, $\partial L/\partial q_k=0$. Una mala elección de coordenadas oculta la simetría.

\vspace{0.4em}
\begin{exampleblock}{En las cuatro clases}
\begin{itemize}
  \item Dos cuerpos: $(\mathbf r_1,\mathbf r_2)\to(\mathbf R,\mathbf r)$; la invariancia traslacional hace cíclica a $\mathbf R$.
  \item Isotropía: movimiento en un plano, coordenadas polares $(r,\theta)$:
  \[L=\tfrac12\mu(\dot r^2+r^2\dot\theta^2)-V(r)\quad\ec{3.20}\]
  \item Oscilaciones pequeñas: $q=q_0+\eta$, con $\eta$ medida desde el equilibrio.
\end{itemize}
\end{exampleblock}
\end{frame}

% ---------------------------------------------------------------------
\begin{frame}{\paso{reducir}{Paso 2} Cantidades conservadas}
\small
\textbf{Movimiento general (Noether).}
\begin{itemize}
  \item $q_k$ cíclica $\Rightarrow$ $p_k=\partial L/\partial\dot q_k=$ cte.
  \item $\partial L/\partial t=0$ $\Rightarrow$ se conserva la función energía; coincide con $E=T+V$ cuando las ligaduras no dependen del tiempo y $V$ no depende de las velocidades.
\end{itemize}
Si hay tantas integrales independientes como grados de libertad, el problema es integrable.

\vspace{0.4em}
\begin{exampleblock}{En las cuatro clases}
\begin{itemize}
  \item Dos cuerpos: $\mathbf P_T$ (3), dirección de $\mathbf l$ (1), $l=\mu r^2\dot\theta$ (1), $E$ (1): seis integrales \ec{3.23, 3.25}.
  \item La 2.ª ley de Kepler es la conservación de $l$:\quad $\dfrac{dA}{dt}=\dfrac{l}{2\mu}$ \ec{3.148}
  \item Perturbación: $l$ se \emph{mantiene fijo}; solo $E$ se eleva por encima de $E_0$.
\end{itemize}
\end{exampleblock}
\end{frame}

% ---------------------------------------------------------------------
\begin{frame}{\paso{reducir}{Paso 3} Reducción a un grado de libertad}
\small
Elimine la velocidad cíclica usando su momento conservado, \emph{en la energía o en la ecuación de movimiento}:
\[
E=\tfrac12 c\,\dot q^2+\Vef(q),\qquad
\Vef(r)=V(r)+\frac{l^2}{2\mu r^2}\quad\ec{3.26}
\]
\begin{columns}[T]
\begin{column}{0.48\textwidth}
\begin{block}{Interpretación}
\footnotesize
El término centrífugo \emph{no} es una ley de fuerza. Es la energía cinética angular, $\tfrac12\mu r^2\dot\theta^2$, reescrita con el $l$ conservado.
\end{block}
\end{column}
\begin{column}{0.48\textwidth}
\begin{alertblock}{Trampa de signo}
\footnotesize
Sustituir $\dot\theta=l/\mu r^2$ en $L$ antes de variar produce $-l^2/2\mu r^2$ en el potencial: el signo equivocado. Elimine en $E$, en la ecuación de movimiento o use el routhiano $R=L-p_\theta\dot\theta$.
\end{alertblock}
\end{column}
\end{columns}
\vspace{0.4em}
{\footnotesize Forma genérica usada en oscilaciones pequeñas: $L=\tfrac12 c\,\dot q^2-\Vef(q)$ \ec{4.1}, donde $c$ reúne los parámetros inerciales.}
\end{frame}

% ---------------------------------------------------------------------
\begin{frame}{\paso{analizar}{Paso 4} Lectura de $\Vef$ y cuadratura}
\small
\begin{columns}[T]
\begin{column}{0.46\textwidth}
\textbf{Lectura cualitativa} (exacta):
\begin{itemize}
  \item Puntos de retorno: raíces de $E=\Vef(q)$.
  \item Movimiento acotado o no acotado según el nivel de $E$.
\end{itemize}
\textbf{Cuadratura} (exacta):
\[
t=\sqrt{\tfrac{\mu}{2}}\int\frac{dr}{\sqrt{E-\Vef(r)}}
\]
\[
\theta=\frac{l}{\sqrt{2\mu}}\int\frac{dr}{r^2\sqrt{E-\Vef(r)}}+\theta_0
\]
{\centering\ec{3.28, 3.33}\par}
\end{column}
\begin{column}{0.52\textwidth}
\begin{exampleblock}{Kepler: $\Vef=-k/r+l^2/2\mu r^2$}
\footnotesize
\begin{itemize}
  \item Cónica $q/r=1+e\cos\theta$, con
  \[e=\sqrt{1+\tfrac{2El^2}{\mu k^2}},\quad q=\tfrac{l^2}{\mu k}\ \ec{3.116}\]
  \item $E<0$: elipse, $a=k/2|E|$ (1.ª ley de Kepler); $E=0$: parábola; $E>0$: hipérbola.
  \item Dependencia temporal: $\omega t=\psi-e\sin\psi$, $r=a(1-e\cos\psi)$ \ec{3.166}.
\end{itemize}
\end{exampleblock}
\end{column}
\end{columns}
\end{frame}

% ---------------------------------------------------------------------
\begin{frame}{\paso{analizar}{Paso 5} Equilibrio y estabilidad}
\small
\[
\Vef'(q_0)=0\ \ \text{(equilibrio)},\qquad \Vef''(q_0)>0\ \ \text{(estable)}\qquad\ec{4.4, 4.5}
\]
\begin{columns}[T]
\begin{column}{0.5\textwidth}
\textbf{Fuerzas centrales.} Una órbita circular es un equilibrio del movimiento \emph{radial}:
\[
\left.\frac{\partial V}{\partial r}\right|_{r_0}=\frac{l^2}{\mu r_0^3}\quad\ec{3.199}
\]
Con la fuerza $f(r)=-V'(r)$, la estabilidad exige
\[
\frac{3f(r_0)}{r_0}+f'(r_0)<0\quad\ec{3.203}
\]
\end{column}
\begin{column}{0.47\textwidth}
\begin{exampleblock}{Fuerza $f=-k\,r^{n}$, $k>0$}
\footnotesize
La condición se reduce a $n>-3$.
\begin{itemize}
  \item Gravedad, $n=-2$: estable.
  \item Resorte isótropo, $n=1$: estable.
  \item $n\le-3$: la barrera centrífuga no retiene la órbita; la partícula cae hacia el centro o escapa en espiral.
\end{itemize}
\end{exampleblock}
\end{column}
\end{columns}
\end{frame}

% ---------------------------------------------------------------------
\begin{frame}{\paso{linealizar}{Paso 6} Linealización}
\small
Tome $q=q_0+\eta$ con $|\eta|$ pequeño y desarrolle hasta segundo orden; el término lineal se anula por el paso 5:
\[
\Vef(q)\approx \Vef(q_0)+\tfrac12 K\eta^2,\ \ K\equiv\Vef''(q_0)
\ \Longrightarrow\
\ddot\eta+\omega^2\eta=0,\ \ \omega^2=\frac{K}{c}\ \ \ec{4.9}
\]
\begin{columns}[T]
\begin{column}{0.48\textwidth}
\begin{exampleblock}{Oscilaciones radiales}
\footnotesize
\[\omega_r^2=\frac{1}{\mu}\left.\frac{\partial^2\Vef}{\partial r^2}\right|_{r_0}\ \ec{3.211}\]
El oscilador armónico es el modelo \emph{local} universal de todo mínimo estable.
\end{exampleblock}
\end{column}
\begin{column}{0.48\textwidth}
\begin{block}{Una sorpresa de Kepler}
\footnotesize
Con $u=1/r$, la órbita de Kepler satisface
\[u''+u=\frac{\mu k}{l^2}\ \ec{3.102}\]
un oscilador armónico en el \emph{ángulo} $\theta$. Es \emph{exacto}, no linealizado; por eso la solución es una cónica.
\end{block}
\end{column}
\end{columns}
\end{frame}

% ---------------------------------------------------------------------
\begin{frame}{Paso 6: varios grados de libertad}
\small
Cerca de una configuración estable $\mathbf q_0$, con $\boldsymbol\eta=\mathbf q-\mathbf q_0$:
\[
L\approx\tfrac12\,\dot{\boldsymbol\eta}^{\mathsf T}\,\mathsf T\,\dot{\boldsymbol\eta}-\tfrac12\,\boldsymbol\eta^{\mathsf T}\,\mathsf V\,\boldsymbol\eta,
\qquad
\mathsf V_{ij}=\left.\frac{\partial^2V}{\partial q_i\partial q_j}\right|_{\mathbf q_0},
\]
donde $\mathsf T_{ij}$ es la matriz cinética evaluada en $\mathbf q_0$. El ensayo $\boldsymbol\eta=\mathbf a\,e^{i\omega t}$ conduce al problema de autovalores generalizado
\[
(\mathsf V-\omega^2\mathsf T)\,\mathbf a=0,\qquad \det(\mathsf V-\omega^2\mathsf T)=0 .
\]
\begin{block}{La misma lógica, en forma matricial}
\footnotesize
\begin{itemize}
  \item Estabilidad $\Leftrightarrow$ $\mathsf V$ definida positiva $\Leftrightarrow$ todas las $\omega_k^2>0$.
  \item Las coordenadas normales desacoplan el movimiento en osciladores independientes.
  \item El resultado unidimensional $\omega^2=K/c$ es el caso $1\times1$.
\end{itemize}
\end{block}
\end{frame}

% ---------------------------------------------------------------------
\begin{frame}{\paso{linealizar}{Paso 7} Comparación de frecuencias}
\small
Dos movimientos periódicos independientes, aquí el radial y el angular, se combinan en una sola órbita. Su \textbf{cociente} decide la geometría:
\begin{gather*}
\Delta\theta=\dot\theta_0\,T_r=2\pi\,\frac{\dot\theta_0}{\omega_r}\quad\ec{3.212}\\
\text{órbita cerrada}\iff\frac{\dot\theta_0}{\omega_r}=\frac{m}{n}\in\mathbb Q\quad\ec{3.216}
\end{gather*}
\begin{columns}[T]
\begin{column}{0.48\textwidth}
\footnotesize
\begin{itemize}
  \item Cociente racional: la órbita se cierra tras $n$ periodos radiales y $m$ revoluciones.
  \item Cociente irracional: órbita cuasiperiódica que nunca se cierra.
  \item Varios modos normales: la misma regla vale para los cocientes $\omega_j/\omega_k$ (figuras de Lissajous).
\end{itemize}
\end{column}
\begin{column}{0.49\textwidth}
\begin{exampleblock}{Escalas de tiempo}
\footnotesize
3.ª ley de Kepler: $T_p=2\pi\sqrt{\mu a^3/k}$ \ec{3.164}.\\[0.2em]
Similitud mecánica para $V\propto r^{\nu}$: $T\propto a^{\,1-\nu/2}$ \ec{3.196}.\\[0.2em]
$\nu=-1$ (gravedad): $T\propto a^{3/2}$; $\nu=2$ (oscilador): $T$ no depende de la amplitud.
\end{exampleblock}
\end{column}
\end{columns}
\end{frame}

% ---------------------------------------------------------------------
\begin{frame}{\paso{validar}{Paso 8} Validación}
\small
\begin{columns}[T]
\begin{column}{0.5\textwidth}
\textbf{Lista para cada resultado}
\footnotesize
\begin{enumerate}
  \item \textbf{Dimensiones} de cada expresión nueva.
  \item \textbf{Signos}: atractivo o repulsivo; término centrífugo positivo.
  \item \textbf{Límites}: perturbación $\to0$, $e\to0$, amplitud $\to0$, $E\to0$.
  \item \textbf{Carácter}: ¿exacto (cuadratura), linealizado o numérico?
  \item \textbf{Consistencia}: ¿coincide el resultado lineal con la integral exacta en su límite?
  \item \textbf{Conservación}: ¿se conservan $E$ y $l$ en toda verificación numérica?
\end{enumerate}
\end{column}
\begin{column}{0.47\textwidth}
\begin{alertblock}{Dominio de validez}
\footnotesize
La linealización vale para $|\eta|$ mucho menor que la escala en la que cambia $\Vef''$. Para amplitud finita, el ángulo apsidal exacto
\[
\Delta\theta=\frac{2l}{\sqrt{2\mu}}\int_{r_{\min}}^{r_{\max}}\frac{dr}{r^2\sqrt{E-\Vef}}\ \ec{3.218}
\]
en general depende de $E$. Un cierre hallado por linealización puede no sobrevivir a amplitud finita.
\end{alertblock}
\end{column}
\end{columns}
\end{frame}

% =====================================================================
\section{El flujo de trabajo en acción}
% =====================================================================

\begin{frame}{Hilo 1: Kepler en los ocho pasos}
\footnotesize
\renewcommand{\arraystretch}{1.25}
\begin{tabular}{@{}L{2.55cm}L{8.6cm}@{}}
\toprule
\textbf{Paso} & \textbf{Resultado para $V(r)=-k/r$}\\
\midrule
\paso{reducir}{1} Simetría & Traslaciones y rotaciones: coordenadas del CM, masa reducida $\mu$, polares planas $(r,\theta)$\\
\paso{reducir}{2} Conservación & $\mathbf P_T$, $\mathbf l$, $E$: seis integrales para seis grados de libertad $\Rightarrow$ integrable\\
\paso{reducir}{3} Reducción & $\Vef(r)=-k/r+l^2/2\mu r^2$ \ec{3.119}\\
\paso{analizar}{4} Lectura de $\Vef$ & Puntos de retorno $r_{\min,\max}=a(1\mp e)$; órbitas cónicas; ecuación de Kepler\\
\paso{analizar}{5} Equilibrio & $r_0=l^2/\mu k$, $E_0=-\mu k^2/2l^2$; $\Vef''(r_0)=k/r_0^3>0$: estable\\
\paso{linealizar}{6} Linealización & $\omega_r^2=k/\mu r_0^3$\\
\paso{linealizar}{7} Frecuencias & $\dot\theta_0^2=l^2/\mu^2r_0^4=k/\mu r_0^3=\omega_r^2$ $\Rightarrow$ $\Delta\theta=2\pi$: \textbf{sin precesión}\\
\paso{validar}{8} Validación & La órbita exacta (paso 4) lo confirma: el perihelio queda fijo para todo $e$, no solo si $e\to0$\\
\bottomrule
\end{tabular}

\vspace{0.6em}
\small
La degeneración $\omega_r=\dot\theta_0$ es la huella dinámica de una cantidad conservada adicional: el vector de Laplace--Runge--Lenz (Sec.~3.7).
\end{frame}

% ---------------------------------------------------------------------
\begin{frame}{Hilo 2: órbitas perturbadas, $f\propto r^{n}$}
\small
\begin{columns}[T]
\begin{column}{0.53\textwidth}
\footnotesize
Fuerza $f(r)=-k\,r^{n}$, $k>0$; órbita circular de radio $r_0$.

\vspace{0.3em}
\paso{analizar}{5} Equilibrio: $l^2=\mu k\,r_0^{\,n+3}$; estable si y solo si $n>-3$.

\vspace{0.3em}
\paso{linealizar}{6} Linealización, usando $\Vef''=-f'-3f/r$:
\[
\omega_r^2=\frac{k\,(n+3)\,r_0^{\,n-1}}{\mu}
\]
\paso{linealizar}{7} Comparación con $\dot\theta_0^2=\dfrac{l^2}{\mu^2r_0^4}=\dfrac{k\,r_0^{\,n-1}}{\mu}$:
\[
\boxed{\ \Delta\theta=\frac{2\pi}{\sqrt{n+3}}\ }
\]
\end{column}
\begin{column}{0.45\textwidth}
\begin{exampleblock}{\paso{validar}{8} Validación}
\footnotesize
\begin{itemize}
  \item $n=-2$ (gravedad): $\Delta\theta=2\pi$. Coincide con el hilo 1.
  \item $n=1$ (oscilador isótropo): $\Delta\theta=\pi$. Elipse centrada: dos ciclos radiales por revolución.
  \item $n\to-3^{+}$: $\Delta\theta\to\infty$, el borde de la estabilidad.
\end{itemize}
\end{exampleblock}
\begin{block}{Conexión con Bertrand}
\footnotesize
El cierre de órbitas casi circulares para todo $r_0$ exige $\sqrt{n+3}$ racional. Los órdenes superiores en la amplitud dejan solo $n=-2$ y $n=1$.
\end{block}
\end{column}
\end{columns}
\end{frame}

% ---------------------------------------------------------------------
\begin{frame}{¿Exacto, aproximado o cualitativo?}
\scriptsize
\renewcommand{\arraystretch}{1.2}
\begin{tabular}{@{}L{4.3cm}L{1.9cm}L{4.9cm}@{}}
\toprule
\textbf{Resultado} & \textbf{Carácter} & \textbf{Dónde falla o cambia}\\
\midrule
Conservación de $l$; 2.ª ley de Kepler & Exacto & Fuerzas no centrales\\
Órbitas cónicas; 1.ª ley de Kepler & Exacto & Todo $V$ no proporcional a $1/r$\\
$T_p=2\pi\sqrt{\mu a^3/k}$; 3.ª ley de Kepler & Exacto & Otros potenciales (la similitud da otro exponente)\\
Puntos de retorno y clasificación de órbitas a partir de $\Vef$ & Exacto, cualitativo & ---\\
$\omega^2=\Vef''(q_0)/c$ & Linealizado & Amplitud finita: correcciones anarmónicas\\
$\Delta\theta=2\pi\,\dot\theta_0/\omega_r$ & Linealizado & Excentricidad finita: usar la Ec.~3.218\\
Frecuencias de los modos normales & Linealizado & Gran amplitud; acoplamiento de modos\\
\bottomrule
\end{tabular}
\vspace{0.6em}

\small
Los pasos 4--5 dan información exacta y global. Los pasos 6--7 dan información aproximada y local. En el paso 8 se decide en cuál se apoya una afirmación.
\end{frame}

% ---------------------------------------------------------------------
\begin{frame}{Errores conceptuales frecuentes}
\small
\begin{itemize}
  \item \paso{reducir}{3} \comillas{El término centrífugo es un potencial real.} Es energía cinética angular reescrita con el $l$ conservado; solo existe para $l$ fijo.
  \item \paso{reducir}{3} Sustituir $p_\theta$ en $L$ antes de variar, lo que invierte el signo de la barrera.
  \item \paso{analizar}{5} Evaluar la estabilidad con $V''$ en lugar de $\Vef''$.
  \item \paso{analizar}{5} Perturbar una órbita circular cambiando $l$ sin advertirlo.
  \item \paso{linealizar}{7} Tomar $\omega_r=\dot\theta_0$ en el problema de Kepler como una coincidencia numérica y no como una simetría.
  \item \paso{linealizar}{7} Creer que las oscilaciones pequeñas requieren un resorte físico: \emph{cualquier} mínimo estable se comporta como uno.
  \item \paso{validar}{8} Suponer que una condición de cierre linealizada vale para cualquier amplitud.
\end{itemize}
\end{frame}

% ---------------------------------------------------------------------
\begin{frame}{Preguntas conceptuales}
\small
Discútalas en parejas; justifique cada respuesta indicando el paso del flujo de trabajo que utilizó.
\begin{enumerate}
  \item Una partícula se mueve en $V(r)=-k/r^4$, $k>0$. ¿Es estable una órbita circular? Responda sin calcular $\Vef''$ explícitamente.
  \item Para cierta fuerza central se obtiene $\omega_r/\dot\theta_0=\sqrt2$. ¿Se cierra una órbita circular ligeramente perturbada? ¿Qué aspecto tiene?
  \item ¿Por qué conocer $\omega^2=\Vef''(q_0)/c$ no basta para saber si se cierran las órbitas de amplitud finita?
  \item La 2.ª ley de Kepler vale para toda fuerza central. ¿Por qué no ocurre lo mismo con la 1.ª y la 3.ª?
  \item ¿Qué paso del flujo de trabajo cambia, y cómo, en un sistema \emph{sin} simetrías continuas?
\end{enumerate}
\end{frame}

% ---------------------------------------------------------------------
\begin{frame}{El paso 8 como hábito}
\small
La lista del paso 8 se aplica a \emph{cualquier} deducción que usted lea: un libro de texto, la solución de un compañero o una respuesta generada por IA.

\vspace{0.3em}
Antes de aceptar un resultado, pregúntese:
\begin{itemize}
  \item ¿Tiene las dimensiones correctas?
  \item ¿Se reduce a un caso conocido (Kepler, oscilador isótropo) en el límite adecuado?
  \item ¿Es exacto o linealizado, y se dice explícitamente?
  \item ¿Se eliminó correctamente el término centrífugo?
\end{itemize}
\vspace{0.3em}
\begin{alertblock}{Hábito por desarrollar}
\footnotesize
Una notación pulida no prueba que el resultado sea correcto. En este tema, una sola verificación de límites detecta la mayoría de los errores de signo y de factores. Por ejemplo, todo $\Delta\theta$ propuesto debe dar $2\pi$ para $V=-k/r$ y $\pi$ para $V\propto r^2$.
\end{alertblock}
\end{frame}

% =====================================================================
\section{Perspectivas y resumen}
% =====================================================================

\begin{frame}{Hacia las próximas unidades}
\footnotesize
\begin{columns}[T]
\begin{column}{0.48\textwidth}
\textbf{Dinámica hamiltoniana (pasos 2--3).} Con $p_\theta=l$ fijo,
\[
H=\frac{p_r^2}{2\mu}+\Vef(r):
\]
la reducción se convierte en movimiento en un plano de fase $(r,p_r)$, y el paso 4 en la lectura de sus curvas de nivel.

\vspace{0.3em}
\textbf{Cuerpos rígidos (pasos 5--6).} La estabilidad de la rotación alrededor de los ejes principales usa la misma lógica de linealizar y evaluar.
\end{column}
\begin{column}{0.48\textwidth}
\textbf{Variables de acción--ángulo (paso 7).} Las frecuencias son $\omega_k=\partial H/\partial J_k$. Para Kepler,
\[
H=-\frac{\mu k^2}{2\,(J_r+J_\theta)^2}
\]
de modo que $\omega_r=\omega_\theta$: la degeneración del hilo 1, ahora explicada estructuralmente (Sec.~6.10).

\vspace{0.3em}
\textbf{Perturbaciones y caos (paso 8).} Los cocientes de frecuencias conmensurables son donde las perturbaciones actúan con más fuerza: el punto de partida de la teoría KAM y del movimiento caótico.
\end{column}
\end{columns}
\end{frame}

% ---------------------------------------------------------------------
\begin{frame}{Resumen}
\small
\begin{columns}[T]
\begin{column}{0.55\textwidth}
\textbf{El flujo de trabajo}
\footnotesize
\begin{enumerate}\setlength{\itemsep}{0pt}
  \item Simetrías y coordenadas adaptadas
  \item Cantidades conservadas (Noether)
  \item Reducción a $\Vef(q)$
  \item Puntos de retorno y cuadratura
  \item Equilibrio y estabilidad
  \item Linealización: $\omega^2=\Vef''/c$
  \item Cocientes de frecuencias: precesión, cierre, periodos
  \item Validación: dimensiones, límites, exactitud
\end{enumerate}
\end{column}
\begin{column}{0.42\textwidth}
\begin{block}{Ecuaciones centrales}
\footnotesize
\[\Vef=V+\frac{l^2}{2\mu r^2}\]
\[\omega_r^2=\frac{1}{\mu}\Vef''(r_0)\]
\[\Delta\theta=2\pi\,\frac{\dot\theta_0}{\omega_r}\]
\end{block}
\end{column}
\end{columns}
\vspace{0.3em}
Las leyes de Kepler, el problema de Kepler, las perturbaciones orbitales y las oscilaciones pequeñas son cuatro aplicaciones de este único método.
\end{frame}

\end{document}
